% theme: stoichiometric_relationships_with_and_without_excess_reactants
\MajorQuestion{化学反応の量的関係}
\SetRowSpace{4mm}

\TextTall{原子量は ${\rm H}=1.0$、${\rm C}=12$、${\rm N}=14$、${\rm O}=16$、${\rm Ca}=40$ とする。また、アボガドロ定数を $6.0\times10^{23}\,{\rm mol}^{-1}$ とし、標準状態における気体 $1\,{\rm mol}$ の体積を $22.4\,{\rm L}$ とする。}

\TypeQuestion{$2\mathrm{H_2}+\mathrm{O_2}\longrightarrow2\mathrm{H_2O}$ について、次の問いに答えなさい。}{hydrogen_combustion_without_excess}
\QQRow{水素 $0.80\,{\rm mol}$ と過不足なく反応する酸素の物質量}{$0.40\,{\rm mol}$}{水素 $4.0\,{\rm g}$ から生じる水の質量}{$36\,{\rm g}$}
\QQRow{標準状態で酸素 $11.2\,{\rm L}$ から生じる水の物質量}{$1.0\,{\rm mol}$}{水素分子 $3.0\times10^{23}$個から生じる水分子の個数}{$3.0\times10^{23}$個}

\TypeQuestion{$\mathrm{CaCO_3}\longrightarrow\mathrm{CaO}+\mathrm{CO_2}$ について、次の問いに答えなさい。}{calcium_carbonate_decomposition}
\QQRow{炭酸カルシウム $0.30\,{\rm mol}$ から生じる二酸化炭素の物質量}{$0.30\,{\rm mol}$}{炭酸カルシウム $25\,{\rm g}$ から生じる酸化カルシウムの質量}{$14\,{\rm g}$}
\QQRow{炭酸カルシウム $20\,{\rm g}$ から生じる二酸化炭素の標準状態での体積}{$4.48\,{\rm L}$}{標準状態で二酸化炭素 $6.72\,{\rm L}$ を生じる炭酸カルシウムの質量}{$30\,{\rm g}$}

\TypeQuestion{$\mathrm{N_2}+3\mathrm{H_2}\longrightarrow2\mathrm{NH_3}$ について、次の問いに答えなさい。}{ammonia_synthesis_without_excess}
\QQRow{窒素 $0.50\,{\rm mol}$ から生じるアンモニアの物質量}{$1.0\,{\rm mol}$}{窒素 $14\,{\rm g}$ から生じるアンモニアの質量}{$17\,{\rm g}$}
\QQRow{標準状態で水素 $6.72\,{\rm L}$ から生じるアンモニアの標準状態での体積}{$4.48\,{\rm L}$}{アンモニア分子 $3.0\times10^{23}$個を生じるのに必要な窒素分子の個数}{$1.5\times10^{23}$個}

\AnswerColumnBreak
\LessonQuestionBreak

\TypeQuestion{$2\mathrm{H_2}+\mathrm{O_2}\longrightarrow2\mathrm{H_2O}$ について、次の問いに答えなさい。}{hydrogen_combustion_with_excess}
\TextTall{水素 $3.0\,{\rm mol}$ と酸素 $1.0\,{\rm mol}$ を混合し、完全に反応させた。}
\QQRow{不足する反応物}{酸素}{反応した水素の物質量}{$2.0\,{\rm mol}$}
\QQRow{生じた水の質量}{$36\,{\rm g}$}{反応後に残った水素の物質量}{$1.0\,{\rm mol}$}

\TypeQuestion{$\mathrm{N_2}+3\mathrm{H_2}\longrightarrow2\mathrm{NH_3}$ について、次の問いに答えなさい。}{ammonia_synthesis_with_excess}
\TextTall{窒素 $28\,{\rm g}$ と水素 $3.0\,{\rm g}$ を混合し、完全に反応させた。}
\QQRow{不足する反応物}{水素}{反応した窒素の質量}{$14\,{\rm g}$}
\QQRow{生じたアンモニアの質量}{$17\,{\rm g}$}{反応後に残った窒素の質量}{$14\,{\rm g}$}

\TypeQuestion{$2\mathrm{CO}+\mathrm{O_2}\longrightarrow2\mathrm{CO_2}$ について、次の問いに答えなさい。}{carbon_monoxide_combustion_with_excess}
\TextTall{標準状態で一酸化炭素 $11.2\,{\rm L}$ と酸素 $8.96\,{\rm L}$ を混合し、完全に反応させた。}
\QQRow{不足する反応物}{一酸化炭素}{反応した酸素の標準状態での体積}{$5.60\,{\rm L}$}
\QQRow{生じた二酸化炭素の標準状態での体積}{$11.2\,{\rm L}$}{反応後に残った酸素の標準状態での体積}{$3.36\,{\rm L}$}
